%% LyX 2.3.6.2 created this file.  For more info, see http://www.lyx.org/.
%% Do not edit unless you really know what you are doing.
\documentclass[english]{article}
\usepackage[T1]{fontenc}
\usepackage[latin9]{inputenc}
\usepackage{bm}
\usepackage{amsmath}
\usepackage{amssymb}
\usepackage{babel}
\begin{document}
To evaluate the incoherent contribution to the conductivity we start
by fourier transforming the current operators that are associated
to the kondo and magnetoelastic assisted hoppings between Pd layers.
These particle-current operators are given by 

\begin{equation}
\boldsymbol{J}^{1}\left(\omega_{n}\right)=\frac{1}{\sqrt{N\beta}}\sum_{\begin{array}{c}
\boldsymbol{k}\boldsymbol{k}'\alpha\beta\\
i\nu_{l},i\nu_{n}
\end{array}}\boldsymbol{f}\left(\boldsymbol{k},\boldsymbol{k}'\right)p_{\boldsymbol{k}\alpha}^{\dagger}\left(i\nu_{l}\right)\left[\bm{S}_{\boldsymbol{k}-\boldsymbol{k}'}\left(i\nu_{l}-i\nu_{n}\right)\cdot\bm{\sigma}_{\alpha\beta}\right]p_{\boldsymbol{k}'\beta}\left(i\nu_{n}+i\omega_{n}\right)
\end{equation}

\begin{align}
\boldsymbol{J}^{2}\left(\omega_{n}\right) & =\frac{1}{N\beta}\sum_{\begin{array}{c}
\boldsymbol{k}\boldsymbol{k}'\boldsymbol{q}\alpha\beta i\\
i\nu_{l},i\mu_{m},i\nu_{n}
\end{array}}\boldsymbol{\Lambda}_{i}\left(\boldsymbol{k},\boldsymbol{k}'\right)X_{i,\boldsymbol{k}-\boldsymbol{k}'-\boldsymbol{q}}\left(i\nu_{l}-i\mu_{m}-i\nu_{n}\right)p_{\boldsymbol{k}\alpha}^{\dagger}\left(i\nu_{l}\right)\left[\bm{S}_{\boldsymbol{q}}\left(i\mu_{m}\right)\cdot\bm{\sigma}_{\alpha\beta}\right]p_{\boldsymbol{k}'\beta}\left(i\nu_{n}+i\omega_{n}\right)
\end{align}

where the bare vertices are given by 
\begin{align}
\boldsymbol{f}\left(\boldsymbol{k},\boldsymbol{k}'\right) & =J_{K}\left(\frac{2a_{z}}{\hbar}\right)\hat{z}\sin\left(k_{z}a_{z}+k_{z}'a_{z}\right),
\end{align}

and 

\begin{align}
\boldsymbol{\Lambda}_{+}\left(\boldsymbol{k},\boldsymbol{k}'\right) & =\boldsymbol{\Lambda}_{-}\left(\boldsymbol{k},\boldsymbol{k}'\right)=-\lambda\boldsymbol{f}\left(\boldsymbol{k},\boldsymbol{k}'\right)
\end{align}

we additionally parametrized the magnetoelastic interaction as $\tilde{\alpha}=J_{k}\lambda$
. Note that we have two phonon operators per local moment since there
are two bond streching modes, one that connects to the upper and one
that connects to the lower Pd layers. With these definitions, we proceed
to calculate the current-current correlation functions 

\begin{equation}
\Pi_{\boldsymbol{k},i}^{zz}\left(i\omega_{n}\right)=\frac{-1}{\beta V}\langle J_{z}^{i}(\tilde{q})J_{z}^{i}(-\tilde{q})\rangle=\Pi_{1}+\Pi_{2}
\end{equation}

There are two diagrams that contibute to leading order, these are
respectively:

where the dotted line corresponds to the $\boldsymbol{J}^{1}$ current,
the dashed line is a spin propagator, the zig-zag line is a $\boldsymbol{J}^{2}$
current, and the wiggly line is a phonon propagator. 

For the first diagram we consider only the Kondo vertex, the diagram
itself is given by 

\begin{equation}
\Pi_{\boldsymbol{k},1}^{zz}\left(i\omega_{n}\right)=\frac{1}{\beta V}\frac{1}{\beta N}\sum_{\boldsymbol{q}\boldsymbol{p}}\sum_{i\nu_{n}iz_{n}}\left|f^{z}\left(\boldsymbol{p}+\boldsymbol{q},\boldsymbol{k}+\boldsymbol{q}\right)\right|^{2}S\left(\boldsymbol{p},iz_{n}\right)\mathcal{G}\left(\boldsymbol{p}+\boldsymbol{q},i\nu_{n}+iz_{n}\right)\mathcal{G}\left(\boldsymbol{k}+\boldsymbol{q},i\omega_{n}+i\nu_{n}\right)
\end{equation}

which comes from the definition of the current from the continutiry
equation. We expect that in the high temperature limit vertex corrections
can be safely ignored in the perturbative regime when quasi-elastic
collisions dominate and there is no momentum structure to the scattering
potentials. In this case we proceed by performing the bosonic and
fermionic matsubara sumations. The result of this is

\begin{align}
\Pi_{\boldsymbol{k},1}^{zz}\left(i\omega_{n}\right) & =\frac{1}{V}\frac{1}{N}\sum_{\boldsymbol{q}\boldsymbol{p}}\left|f^{z}\left(\boldsymbol{p}+\boldsymbol{q},\boldsymbol{k}+\boldsymbol{q}\right)\right|^{2}\int\frac{d\epsilon}{\pi}\int\frac{d\nu}{\pi}\int\frac{dz}{\pi}\chi''\left(\boldsymbol{p},z\right)A\left(\boldsymbol{p}+\boldsymbol{q},\nu\right)A\left(\boldsymbol{k}+\boldsymbol{q},\epsilon\right)\\
 & \ \ \ \ \ \ \ \ \ \ \ \ \ \ \ \ \ \ \ \ \times\left(\frac{n_{F}\left(\epsilon\right)-n_{F}\left(\nu-z\right)}{i\omega_{n}-\epsilon+\nu-z}\right)\left(n_{B}\left(z\right)+n_{F}\left(\nu\right)\right).
\end{align}

We see that overall the current correlation function is a convolution
of electron spectral functions with the spin susceptibility. Once
we perform analytic continuation and take the zero frequency limit,
we get 
\begin{equation}
\sigma_{11}^{c}=\frac{\hbar e^{2}}{VN}\sum_{\boldsymbol{q}\boldsymbol{p}}\left|f^{z}\left(p_{z},q_{z}\right)\right|^{2}\int\frac{d\epsilon}{\pi}\left[\int\frac{dz}{\pi}\chi''\left(\epsilon-z\right)A\left(\boldsymbol{p},z\right)\left(n_{B}\left(\epsilon-z\right)+n_{F}\left(-z\right)\right)\right]A\left(\boldsymbol{q},\epsilon\right)\left(-\frac{\partial n_{F}\left(\epsilon\right)}{\partial\epsilon}\right),
\end{equation}

where we already took the $\boldsymbol{k}\rightarrow0$ limit and
we used the local approximation on $\chi''$ to simplify the momentum
sums. Since the electron spectral functions satisfy $A\left(\boldsymbol{p},\omega\right)\approx A\left(\left(p_{x},p_{y}\right),\omega\right)$
we can pull the momentum sums over $q_{x,y}$ and $p_{x,y}$ past
the bare vertex, and to leading order we have 

\begin{equation}
\sigma_{11}^{c}=\frac{2a_{z}}{a_{pl}^{2}}\hbar e^{2}\int\frac{dp_{z}}{2\pi}\frac{dq_{z}}{2\pi}\left|f^{Z}\left(p_{z},q_{z}\right)\right|^{2}\int\frac{d\epsilon}{\pi}\left[\int\frac{dz}{\pi}\chi''\left(\epsilon-z\right)\rho\left(z\right)\left(n_{B}\left(\epsilon-z\right)+n_{F}\left(-z\right)\right)\right]\rho\left(\epsilon\right)\left(-\frac{\partial n_{F}\left(\epsilon\right)}{\partial\epsilon}\right)
\end{equation}

using 

\[
\sigma_{11}^{c}=\left(\frac{2a_{z}}{a_{pl}^{2}}\frac{e^{2}}{\hbar}\right)J_{K}^{2}\int\frac{d\epsilon}{\pi}\rho\left(\epsilon\right)\left(-\frac{\partial n_{F}\left(\epsilon\right)}{\partial\epsilon}\right)\int\frac{dz}{\pi}\rho\left(\epsilon-z\right)\chi''\left(z\right)\left(n_{B}\left(z\right)+n_{F}\left(z-\epsilon\right)\right)
\]

To evaluate this integral we use the density of states evaluated from
the dispersion in the low energy theory of XXX

Now we focus on the magnetoelastic current correlation functions

\[
\Pi_{\boldsymbol{q},2}^{zz}\left(i\omega_{n}\right)=-\frac{1}{\beta^{3}VN^{2}}\sum_{\begin{array}{c}
\boldsymbol{k}'\boldsymbol{k}\boldsymbol{p}\\
ik_{n}'ik_{n}ip_{n}\\
\ell
\end{array}}\left|\Lambda_{\ell}^{\alpha}\left(\boldsymbol{k},\boldsymbol{k}'+\boldsymbol{q}\right)\right|^{2}D\left(\boldsymbol{k}-\boldsymbol{k}'-\boldsymbol{p},ik_{n}-ip_{n}-ik_{n}'\right)S\left(\boldsymbol{p},ip_{n}\right)\mathcal{G}\left(\boldsymbol{k},ik_{n}\right)\mathcal{G}\left(\boldsymbol{k}'+\boldsymbol{q},ik_{n}'+iq_{n}\right)
\]

where we defined 

\begin{align}
\sigma_{22}^{c} & =-\frac{e^{2}}{VN}\sum_{\boldsymbol{q}\boldsymbol{k}\ell}\int\frac{d\epsilon}{\pi}\frac{d\mu}{\pi}\frac{dz}{\pi}\left|\Lambda_{\ell}^{\alpha}\left(\boldsymbol{q},\boldsymbol{k}\right)\right|^{2}B\left(z\right)\chi''\left(\mu\right)A\left(\boldsymbol{q},\epsilon+z+\mu\right)A\left(\boldsymbol{k},\epsilon\right)\times\\
 & \ \ \ \left[n_{B}\left(-z\right)-n_{B}\left(\mu\right)\right]\left[n_{B}\left(z+\mu\right)+n_{F}\left(\epsilon+z+\mu\right)\right]\left(-\frac{\partial n_{F}\left(\epsilon\right)}{\partial\epsilon}\right)
\end{align}

\begin{align}
\sigma_{22}^{c} & =\left(\frac{2a_{z}}{a_{pl}^{2}}\frac{e^{2}}{\hbar}\right)\left(\frac{\hbar\lambda^{2}}{M\omega_{0}}\right)\left[2J_{K}\right]^{2}\text{csch}\left(\frac{\beta\omega_{0}}{2}\right)\int\frac{d\epsilon}{\pi}\frac{d\mu}{\pi}\rho\left(\epsilon\right)\cosh\left(\frac{\beta\epsilon}{2}\right)\left(-\frac{\partial n_{F}\left(\epsilon\right)}{\partial\epsilon}\right)\times\\
 & S\left(\mu\right)\left[\rho\left(\epsilon+\omega_{0}+\mu\right)\text{sech}\left(\beta\frac{\epsilon+\mu+\omega_{0}}{2}\right)+\rho\left(\epsilon-\omega_{0}+\mu\right)\text{sech}\left(\beta\frac{\epsilon+\mu-\omega_{0}}{2}\right)\right]\left[\frac{\beta\mu}{2}\text{csch}\left(\frac{\beta\mu}{2}\right)\right]
\end{align}

\end{document}
